% LCA in O(1) with RMQ

Lowest Common Ancestor (LCA) queries on a tree in $O(1)$ with $O(N \log N)$ preprocessing using Euler Tour and Sparse Table.
\begin{lstlisting}
struct LCA_O1 {
    vector<int> euler, first, depth;
    vector<vector<int>> st;
    int n, K;

    LCA_O1(int n, int root, const vector<vector<int>>& adj) : n(n) {
        first.resize(n + 1);
        depth.resize(n + 1);
        euler.reserve(2 * n);
        dfs(root, 0, 0, adj);
        
        int m = euler.size();
        K = __lg(m);
        st.assign(K + 1, vector<int>(m));
        for (int i = 0; i < m; i++) st[0][i] = euler[i];
        
        for (int j = 1; j <= K; j++) {
            for (int i = 0; i + (1 << j) <= m; i++) {
                int u = st[j - 1][i], v = st[j - 1][i + (1 << (j - 1))];
                st[j][i] = (depth[u] < depth[v]) ? u : v;
            }
        }
    }

    void dfs(int u, int p, int d, const vector<vector<int>>& adj) {
        first[u] = euler.size();
        euler.push_back(u);
        depth[u] = d;
        for (int v : adj[u]) {
            if (v != p) {
                dfs(v, u, d + 1, adj);
                euler.push_back(u);
            }
        }
    }

    int query(int u, int v) {
        int l = first[u], r = first[v];
        if (l > r) swap(l, r);
        int j = __lg(r - l + 1);
        int node1 = st[j][l], node2 = st[j][r - (1 << j) + 1];
        return (depth[node1] < depth[node2]) ? node1 : node2;
    }

    int get_dist(int u, int v) {
        return depth[u] + depth[v] - 2 * depth[query(u, v)];
    }
};
\end{lstlisting}
